% from main.tex
%%%%%%%%%%  scan idx 614  |  folio 601  %%%%%%%%%%
Soient \(Y_k=\bigcup_{i\in I(k)}X_i\) les composantes connexes de \(X\)~; alors \(f\) se décompose, \(f=\coprod_k f_k:\coprod_k(\coprod_{i\in I(k)}X_i)\longrightarrow Y_k\), et l'assertion en résulte. Enfin, une récurrence immédiate nous ramène au cas d'un morphisme \(X_1\amalg X_2\to X\). \(\sgacircled{1}\) et \(\sgacircled{3}\) résultent alors, quitte à restreindre \(S\) d'après (EGA III 7.8.6), de la proposition suivante.

\textbf{Proposition 2.3.} Soient \(S\) un schéma, \(g:Y\to S\) un morphisme propre et \(Y_1,Y_2\) deux sous-schémas fermés définis par les idéaux \(I_1,I_2\). On suppose que
\begin{enumerate}[label=\alph*)]
\item \(Y,Y_1,Y_2,Y_1\cap Y_2\) sont fidèlement plats et de présentation finie sur \(S\).
\item \(g_*\mathcal O_Y=\mathcal O_S\) et \((g_i)_*\mathcal O_{Y_i}=\mathcal O_S\) universellement (où \(g_i=g|_{Y_i}\), \(i=1,2\)).
\item \(I_1\cap I_2=0\).
\end{enumerate}
Alors \(\Pic_{Y/S}\longrightarrow \Pic_{(Y_1\amalg Y_2)/S}\) est représentable par un morphisme affine de présentation finie.

En effet notons d'abord que les hypothèses sont stables pour tout changement de base. Pour (c), celà résulte des deux suites exactes~:
\[
  0\longrightarrow I_1+I_2\longrightarrow\mathcal O_Y
  \longrightarrow\mathcal O_{Y_1\cap Y_2}\longrightarrow0
\]
et
\[
  0\longrightarrow I_1\cap I_2\longrightarrow I_1\oplus I_2
  \longrightarrow I_1+I_2\longrightarrow0~;
\]
la première montre que \(I_1+I_2\) est plat sur \(S\)~; ensuite la seconde que \(I_1\cap I_2\) est stable par changement de base.

L'assertion est locale pour la topologie f.p.p.f. En faisant le changement de base \(Y_1\cap Y_2\to S\), on peut donc supposer que \(Y_1\cap Y_2/S\) possède une section \(\alpha\), qui définit des sections \(\alpha_i\) de \(Y_i\) et \(\beta\) de \(Y\). Après un changement de base convenable, on est ramené à voir que tout point \(\mu:S\to\Pic_{(Y_1\amalg Y_2)/S}\) a comme image réciproque un sous-foncteur de \(\Pic_{Y/S}\) qui est
